\section{国际商事仲裁}
\begin{enumerate}
  \item 跟国内民商事争议解决方式，没有本质的不同。协商、调解、\zd{仲裁和诉讼}。
  \item 仲裁地。仲裁庭实际开庭审理并作出仲裁裁决的地点。\zd{法院地之于诉讼}。
  \item 被承认执行的可能性远大于法院判决，是一个比诉讼好很多倍的争议解决方式。
  \item 人身密切相关不能申请仲裁。
  \item 仲裁协议效力的认定。
        \begin{enumerate}
          \item \zd{仲裁机构和法院}都有权认定。
          \item 一方仲裁，一方法院 \yjt \zd{法院}裁定。
          \item 管辖法院。仲裁机构所在地、仲裁协议签订地、申请人住所地、被申请人住所地的\zd{中院}。
          \item 时间。仲裁庭\zd{首次开庭前}，方式。贸仲要求\zd{书面}。
          \item 适用法律问题。
                \begin{enumerate}
                  \item 意思自治 \yjt 仲裁机构所在地/仲裁地 \yjt 中国法。
                  \item 意思自治一定要在仲裁条款的\zd{明确表示}。\al{使用 xx 国法进行仲裁}。
                  \item 都认定为无效 + 无效才是无效。有效中院可\zd{直接做出}，无效需\zd{内部报核}直到最高。
                \end{enumerate}
        \end{enumerate}
\end{enumerate}
